\documentclass[10pt,a4paper]{article}
\usepackage[utf8]{inputenc}
\usepackage[T1]{fontenc}
\usepackage[margin=1.8cm,top=2.6cm,bottom=2.2cm]{geometry}
\usepackage{parskip}
\usepackage{microtype}
\usepackage{amsmath,amssymb,mathtools}
\usepackage{array,booktabs,tabularx,multirow}
\usepackage[table]{xcolor}
\usepackage{graphicx,float}
\usepackage{titlesec}
\usepackage{enumitem}
\usepackage{fancyhdr}

\renewcommand\thesection{\Roman{section}}
\numberwithin{equation}{section}

\titleformat{\section}{\large\bfseries}{\thesection.\;}{0em}{}[\titlerule]
\titleformat{\subsection}{\normalsize\bfseries}{}{0em}{}
\titleformat{\subsubsection}{\normalsize\itshape}{}{0em}{}
\titlespacing{\section}{0pt}{1.6ex plus .2ex}{0.6ex}
\titlespacing{\subsection}{0pt}{1.0ex}{0.3ex}

\pagestyle{fancy}
\fancyhf{}
\setlength{\headheight}{15pt}
\fancyhead[L]{\small\textit{Layer~13: Ratiometric Signal and Enzyme Quantification}}
\fancyhead[R]{\small Aaran Nihalani $\bullet$ Eton iGEM Dry Lab}
\fancyfoot[C]{\small\thepage}
\renewcommand{\headrulewidth}{0.35pt}

\newcommand{\cit}[1]{[#1]}
\newcommand{\Eo}{E^{\circ\prime}}
\newcommand{\kB}{k_{\text{B}}}
\newcommand{\kET}{k_{\text{ET}}^{0}}
\newcommand{\Esw}{E_{\text{sw}}}
\newcommand{\Aarea}{A_{\text{area}}}
\newcommand{\AMB}{A_{\text{MB}}}
\newcommand{\AFc}{A_{\text{Fc}}}
\newcommand{\Gtot}{\Gamma_{\text{intact}}}
\newcommand{\kcat}{k_{\text{cat}}}
\newcommand{\Km}{K_{m,\text{surf}}^{\Gamma}}
\newcommand{\kk}{\kappa_k}
\newcommand{\Fk}{F_k}

\begin{document}

%=======================================================================
\begin{center}
{\LARGE\bfseries Layer~13: Ratiometric Signal Processing and Enzyme Quantification}\\[0.4em]
{\large From Peak Areas to Enzyme Concentration with Propagated Uncertainty}\\[0.25em]
{\small\color{gray} Eton iGEM 2026, Dry Lab \quad$|$\quad
Layer~13 of 17-Layer Biosensor Pipeline \quad$|$\quad
\texttt{ratiometric.py}}
\end{center}
\vspace{0.4em}\hrule\vspace{0.7em}

\noindent
Layer~12 delivers two peak areas per SWV scan: $\AMB$ (from the MB redox
label covalently attached to the peptide probe SGKIPRTLTA) and $\AFc$ (from the
ferrocenemethanol internal reference). Layer~13 converts these two numbers into a
single enzyme concentration estimate $[\text{E}]$ with a full uncertainty budget.

The challenge this layer must solve is that the raw peak areas are corrupted
by at least four distinct artefacts: variation in electrode area $A_\text{true}$
between cartridge wells, gain drift in the transimpedance amplifier,
variation in the SWV voltage amplitude $\Delta E_\text{sw}$ over time, and
electrode fouling (non-specific adsorption of plasma proteins that blocks
the electrode surface). The ratiometric scheme described here cancels three
of these artefacts exactly and the fourth approximately, leaving a signal
whose remaining uncertainty comes almost entirely from Poisson-like shot
noise on the peak area fit, which can be propagated analytically.

%=======================================================================
\section{The Raw Ratio $R$: Cancelling Electrode-to-Electrode Variation}
\label{sec:ratio}

\subsection{Term-by-term derivation}

From Layer~11 (Eq.~11.8), the measured MB peak area on electrode $i$ is:
\begin{equation}
  \AMB^{(i)} = n_e^2 \frac{F^2}{2RT}\,f(\kappa_{\text{SW}}^{(i)})\,
  A_\text{true}^{(i)}\,\Gtot^{(i)}\,\Esw^{(i)}
  \label{eq:AMB}
\end{equation}
where $n_e = 2$, $f(\kappa_\text{SW})$ is the dimensionless kinetic factor
(between 0 and~1, from Layer~11, Section~III), $A_\text{true}^{(i)}$ is the
actual gold area of electrode $i$, $\Gtot^{(i)}$ is the surviving intact
probe surface coverage, and $\Esw^{(i)}$ is the SWV amplitude delivered
by the potentiostat at the time of the scan.

The ferrocenemethanol reference peak area, from the same electrode and the
same scan, obeys the same equation with $n_e = 1$, its own formal potential,
its own kinetic factor, its own tether length (set by design to give
$f(\kappa_\text{SW}^\text{Fc}) \approx 1$ at the chosen $f_\text{SWV}$),
and its own surface coverage $\Gamma_\text{Fc}^{(i)}$:
\begin{equation}
  \AFc^{(i)} = 1 \cdot \frac{F^2}{2RT}\,f(\kappa_{\text{SW}}^{\text{Fc},(i)})\,
  A_\text{true}^{(i)}\,\Gamma_\text{Fc}^{(i)}\,\Esw^{(i)}
  \label{eq:AFc}
\end{equation}
Notice that $A_\text{true}^{(i)}$ and $\Esw^{(i)}$ appear identically in
both Eq.~\eqref{eq:AMB} and Eq.~\eqref{eq:AFc}. The ratio:
\begin{equation}
  \boxed{
  R^{(i)} = \frac{\AMB^{(i)}}{\AFc^{(i)}}
  = n_e^2\,\frac{f(\kappa_\text{SW}^{(i)})}{f(\kappa_\text{SW}^{\text{Fc},(i)})}\,
  \frac{\Gtot^{(i)}}{\Gamma_\text{Fc}^{(i)}}
  }
  \label{eq:R}
\end{equation}
The electrode area $A_\text{true}^{(i)}$ and the SWV amplitude $\Esw^{(i)}$
cancel exactly. Physical constants $F, R, T$ also cancel, removing the need
for temperature calibration in this ratio (though Layer~10 still needs $T$
for the $\Eo_\text{MB}(\text{pH})$ prediction).

\subsection{What $R$ still depends on}

The remaining factors in Eq.~\eqref{eq:R} are:
(1) The kinetic ratio $f(\kappa_\text{SW})/f(\kappa_\text{SW}^\text{Fc})$.
Because the Fc tether is chosen to put Fc in the fast-kinetics limit
($\kappa_\text{SW}^\text{Fc} \gg 1$ so $f \approx 1$), and because the
MB kinetic factor $f(\kappa_\text{SW})$ is a function of $\kET \times
t_\text{pulse}$ that is the same for all electrodes of the same design
(since $\kET$ depends on tether length, not on $A_\text{true}$), this ratio
is a device-level constant that cancels when we subsequently take $R(T)/R(0)$
(Section~\ref{sec:rho}).

(2) The coverage ratio $\Gtot/\Gamma_\text{Fc}$.
The Fc SAM and the aptamer-MB SAM are co-deposited, so their surface densities
are correlated: electrodes with more Fc also have more MB. If the Fc coverage
is stable (Fc molecules do not desorb during the assay, as confirmed by
Layer~12 diagnostics), then $\Gamma_\text{Fc}$ is a constant for a given
electrode and the ratio $\Gtot/\Gamma_\text{Fc}$ evolves purely due to
changes in $\Gtot$ (i.e., cleavage or folding of the aptamer probe).

White and Plaxco~\cit{1} confirmed experimentally that taking this ratio
eliminates electrode-to-electrode amplitude variation by over 10-fold
in a chip with deliberate variation in electrode size, validating the
$A_\text{true}$ cancellation.

\subsection{Uncertainty on $R$}

The uncertainty on $R$ from the peak fitting uncertainties is propagated
using the standard formula for the ratio of two independent random variables:
\begin{equation}
  \left(\frac{\sigma_R}{R}\right)^2 =
  \left(\frac{\sigma_{\AMB}}{\AMB}\right)^2 +
  \left(\frac{\sigma_{\AFc}}{\AFc}\right)^2
  \label{eq:sigR}
\end{equation}
This assumes that the MB and Fc areas are independent; they are indeed fitted
from different non-overlapping regions of the voltammogram (Section~12.V),
so their fit errors are uncorrelated. The relative uncertainty on $R$ is
therefore quadrature-added from the relative uncertainties on each area.

Numerically, with $\AMB \approx 0.3$~nA\,V and $\sigma_{\AMB} \approx 3\%$
of $\AMB$ (Layer~12, Eq.~12.12), and with similar relative precision on
$\AFc$: $\sigma_R/R \approx 4.2\%$ per scan.

%=======================================================================
\section{The Normalised Signal $\rho$: Cancelling Probe Density Variation}
\label{sec:rho}

\subsection{Motivation}

Even after forming $R$ (Eq.~\ref{eq:R}), the coverage ratio
$\Gtot/\Gamma_\text{Fc}$ varies between electrodes (different deposition
conditions, different self-assembly kinetics) by up to $\pm 30\%$ (Lubin
and Plaxco~\cit{2}). This means $R(0)$ (the ratio measured before enzyme
exposure) is electrode-specific. To cancel this, we normalise each
electrode's $R(T)$ at time $T$ by its own baseline $R(0)$:
\begin{equation}
  \boxed{
  \rho^{(i)}(T) = \frac{R^{(i)}(T)}{R^{(i)}(0)}
  = \frac{\Gtot^{(i)}(T)}{\Gtot^{(i)}(0)}
  }
  \label{eq:rho}
\end{equation}
The kinetic ratio and Fc coverage cancel in Eq.~\eqref{eq:rho} because they
are constant for a given electrode: $f(\kappa_\text{SW})$ does not change
during the assay (the scan parameters are fixed), and $\Gamma_\text{Fc}$
is stable over the 30-minute assay time (Layer~12 diagnostics confirm
$<2\%$ Fc amplitude drift per scan).

The physical meaning of $\rho$ is exactly the fraction of aptamer probes
that remain intact at time $T$: $\rho = 1$ means no cleavage, $\rho = 0$
means complete cleavage. For the biosensor, we expect $\rho \in [0.1, 1]$
over the clinically relevant enzyme concentration range.

\subsection{Concrete fouling example}

Consider an electrode with 20\% non-specific protein adsorption (fouling)
that covers a fraction of the gold surface. Fouling reduces both $\AMB$
and $\AFc$ by the same fractional amount $f_\text{foul}$:
\begin{align}
  \AMB(T) &= \AMB^\text{clean}(T) \times (1 - f_\text{foul})
  \label{eq:foul_MB}\\
  \AFc(T) &= \AFc^\text{clean}(T) \times (1 - f_\text{foul})
  \label{eq:foul_Fc}
\end{align}
When we form $R = \AMB/\AFc$, the fouling factor $(1-f_\text{foul})$
cancels because it appears in both numerator and denominator.
Similarly, the baseline measurement $R(0)$ also suffered the same
fouling (adsorption is a fast process that saturates before the first
scan). Therefore $\rho = R(T)/R(0)$ also cancels the fouling
factor: the ratio of the ratios is the ratio of the cleavage fractions
with no fouling term surviving.

This cancellation is \emph{exact} only if the fouling is the same on the Fc
and MB regions of the electrode, which is true if the SAM is uniform.
A non-uniform fouling pattern (e.g., a protein adsorbing preferentially
to the aptamer-MB SAM but not to the Fc SAM because of electrostatic
differences) would cause only partial cancellation. Section~\ref{sec:limits}
addresses this limitation.

\subsection{Uncertainty on $\rho$}

\begin{equation}
  \left(\frac{\sigma_\rho}{\rho}\right)^2 =
  \left(\frac{\sigma_{R(T)}}{R(T)}\right)^2 +
  \left(\frac{\sigma_{R(0)}}{R(0)}\right)^2
  \label{eq:sigrho}
\end{equation}
The baseline $R(0)$ is the average of three scans taken before enzyme
addition; averaging $n_0 = 3$ scans reduces $\sigma_{R(0)}$ by
$1/\sqrt{3}$. With $\sigma_R/R \approx 4.2\%$ per scan and three baseline
scans:
\begin{equation}
  \frac{\sigma_{R(0)}}{R(0)} = \frac{4.2\%}{\sqrt{3}} = 2.4\%
  \label{eq:sigrho0}
\end{equation}
Combining with a single on-assay scan:
\begin{equation}
  \frac{\sigma_\rho}{\rho} = \sqrt{(4.2\%)^2 + (2.4\%)^2} = 4.8\%
  \label{eq:sigrho_val}
\end{equation}
This is the irreducible per-electrode fractional uncertainty on $\rho$.
Taking the average over $N_\text{active}$ active electrodes on a chip
reduces it further by $1/\sqrt{N_\text{active}}$.

%=======================================================================
\section{The Double Ratio $F_k$: Cancelling Cleavage Independent of Enzyme}
\label{sec:Fk}

\subsection{Background cleavage and the control electrode}

Even without target enzyme, the aptamer probe is not perfectly stable.
Background cleavage can occur from (a) dissolved nucleases in the sample
matrix (plasma contains substantial nuclease activity), (b) hydrolysis of
the DNA backbone at the assay temperature, and (c) reactive oxygen species
in the wound fluid.

This background cleavage reduces $\rho$ on the active electrode by an
amount $\Delta\rho_\text{bg}$ that is unrelated to the target enzyme.
Unless accounted for, it would overestimate $F_k$ and therefore
overestimate $[\text{E}]$.

The solution is a \emph{control electrode}: an electrode bearing the
D-amino acid enantiomer of the peptide substrate (D-SGKIPRTLTA) in an
otherwise identical SAM. MMP-9, like all natural metalloproteinases, is
stereospecific for L-amino acid substrates: the active-site geometry that
positions the scissile bond over the catalytic zinc ion requires the
L-configuration at each $\alpha$-carbon. The D-enantiomer substrate cannot
adopt this geometry and is therefore not cleaved by MMP-9 (or any natural
L-specific protease in wound fluid) regardless of enzyme concentration.
Non-stereospecific background processes (non-specific nucleases acting on
the thiol linker, gold--thiol bond hydrolysis, and reactive oxygen species)
act on both the L-peptide active electrode and the D-peptide control equally.
The control electrode therefore reports $\rho_\text{control}(T)$ as the
surviving fraction under background conditions alone.

\subsection{Derivation of $F_k$}

The enzyme-specific fractional cleavage is:
\begin{equation}
  \boxed{
  \Fk = 1 - \frac{\rho_\text{active}(T)}{\rho_\text{control}(T)}
  }
  \label{eq:Fk}
\end{equation}
To understand what cancels in Eq.~\eqref{eq:Fk}, write the surviving fractions
explicitly. Let $p_\text{bg}$ be the fraction of probes cleaved by background
mechanisms, and $p_\text{E}$ be the additional fraction cleaved by the target
enzyme. Then:
\begin{align}
  \rho_\text{active} &= (1 - p_\text{bg})(1 - p_\text{E})
  \label{eq:rho_active}\\
  \rho_\text{control} &= (1 - p_\text{bg})
  \label{eq:rho_control}
\end{align}
Substituting into Eq.~\eqref{eq:Fk}:
\begin{equation}
  \Fk = 1 - \frac{(1-p_\text{bg})(1-p_\text{E})}{(1-p_\text{bg})}
  = 1 - (1 - p_\text{E}) = p_\text{E}
  \label{eq:Fk_simp}
\end{equation}
The factor $(1-p_\text{bg})$ cancels exactly, leaving $\Fk = p_\text{E}$:
the fraction of active probes cleaved by the target enzyme alone, regardless
of what background cleavage occurred. This is the quantity that directly
constrains $[\text{E}]$ through the kinetic model.

Lubin and Plaxco~\cit{2} validate this background cancellation experimentally
by showing that control electrodes in the same plasma sample track active
electrodes when no target is added, confirming that the ratio
$\rho_\text{active}/\rho_\text{control}$ is stable at $\approx 1$ in
the absence of enzyme.

\subsection{What does not cancel in $F_k$}

Two sources of artefact are not cancelled by Eq.~\eqref{eq:Fk}:

(1) \emph{Differential fouling.} If the active and control SAMs foul
differently (Section~\ref{sec:rho}), the Fc amplitude on the two electrodes
drifts at different rates. The ratio $\rho_\text{active}/\rho_\text{control}$
then drifts even without enzyme. This is controlled by monitoring the Fc
peak amplitude on both electrodes; a difference $>10\%$ flags differential
fouling and the control correction is invalidated.

(2) \emph{Non-stereospecific metalloprotease activity.} Although natural
MMPs are L-stereospecific, some bacterial proteases and heat-denatured
MMPs may exhibit reduced stereospecificity. Any MMP-9 activity on the
D-enantiomer would cause $\rho_\text{control}$ to decrease with enzyme
concentration, causing Eq.~\eqref{eq:Fk_simp} to underestimate $p_\text{E}$.
Validation requires adding saturating MMP-9 to a chip bearing only
D-enantiomer control electrodes and confirming $\rho_\text{control}$
drops by less than 5\%.

\subsection{Uncertainty on $F_k$}

From Eq.~\eqref{eq:Fk}:
\begin{equation}
  \sigma_{\Fk}^2 = \left(\frac{1}{\rho_\text{control}}\right)^2 \sigma_{\rho_\text{active}}^2
  + \left(\frac{\rho_\text{active}}{\rho_\text{control}^2}\right)^2 \sigma_{\rho_\text{control}}^2
  \label{eq:sigFk_raw}
\end{equation}
If active and control electrodes have equal per-electrode uncertainty
$\sigma_\rho / \rho \approx 4.8\%$ (Eq.~\ref{eq:sigrho_val}) and if
$\rho_\text{active} \approx \rho_\text{control} \approx \rho$ (early time point,
small cleavage), then:
\begin{equation}
  \sigma_{\Fk} \approx (1 - \Fk)\,\sqrt{2}\,
  \frac{\sigma_\rho}{\rho}
  = (1 - \Fk) \times 6.8\%
  \label{eq:sigFk}
\end{equation}
This expression contains the factor $(1-\Fk)$, which decreases toward zero
as $\Fk \to 1$ (complete cleavage). It might seem that the uncertainty also
vanishes, but this is misleading: the fractional uncertainty $\sigma_{\Fk} /
(1 - \Fk)$ is actually $\sqrt{2} \times \sigma_\rho/\rho = 6.8\%$,
constant regardless of $\Fk$. The practical limit of detection is set by
whether $\Fk$ exceeds $3\sigma_{\Fk}$ at zero enzyme concentration, i.e.,
whether $3 \times 6.8\% \approx 20\%$ cleavage can be reliably distinguished
from baseline noise.

%=======================================================================
\section{Inversion to Enzyme Concentration}
\label{sec:inversion}

\subsection{Connecting $F_k$ to the kinetic model}

From Layer~9 (Section~IV), the integrated MM equation gives the fraction
of probes cleaved $F_k = p_\text{E}$ after incubation time $T$:
\begin{equation}
  \kcat[\text{E}]\,T = \Km\ln\!\left(\frac{1}{1-\Fk}\right) + \Gtot(0)\,\Fk
  \label{eq:intMM}
\end{equation}
where $\kcat$ is the catalytic rate constant, $[\text{E}]$ is the enzyme
molar concentration, $T$ is the incubation time, $\Km$ is the surface
Michaelis constant, and $\Gtot(0)$ is the initial probe density.

\subsection{Two limiting regimes and the apparent rate}

In the unsaturated regime ($\Fk \ll 1$, i.e., early time or low enzyme),
Eq.~\eqref{eq:intMM} simplifies. Using $\ln(1/(1-u)) \approx u + u^2/2$
for small $u$:
\begin{equation}
  \kcat[\text{E}]\,T \approx \left(\Km + \Gtot(0)\right)\Fk
  \quad \Rightarrow \quad
  \Fk \approx \frac{\kcat[\text{E}]\,T}{\Km + \Gtot(0)}
  \label{eq:linlimit}
\end{equation}
In this regime, $\Fk$ is linear in $[\text{E}] \times T$, making
calibration straightforward: the slope is $\kcat / (\Km + \Gtot(0))$.

In the saturated regime ($\Gtot \ll \Km$, i.e., low probe density), the
linear term dominates:
\begin{equation}
  \kcat[\text{E}]\,T \approx \Km \Fk
  \quad \Rightarrow \quad
  \Fk \approx \frac{\kcat[\text{E}]\,T}{\Km}
  \label{eq:kmlimit}
\end{equation}
In both limits, we can define an \emph{apparent rate} that depends only on
observables:
\begin{equation}
  \boxed{
  \kk = -\frac{\ln(1-\Fk)}{T}
  }
  \label{eq:kappa}
\end{equation}
$\kk$ has units of $\text{s}^{-1}$ (or $\text{min}^{-1}$ for $T$ in
minutes). It is the rate of the monoexponential process that would
produce the same $\Fk$ in time $T$. In the unsaturated regime where
$\Fk \ll 1$, $\ln(1/(1-\Fk)) \approx \Fk$, so $\kk \approx \Fk/T =
\kcat[\text{E}]/(\Km + \Gtot(0))$, which is linear in $[\text{E}]$.
More generally, $\kk$ is a monotone increasing function of $\Fk$ and
therefore of $[\text{E}]$: the calibration curve $[\text{E}]$ vs $\kk$
is monotone and invertible.

Equation~\eqref{eq:kappa} was chosen as the primary reported observable
rather than $\Fk$ itself for two reasons. First, if the incubation time
$T$ varies between runs (e.g., because the clinician adds sample at a
different time), $\kk$ normalises by $T$ so the calibration curve is
time-independent. Second, Lubin and Plaxco~\cit{2} and Baker~2006~\cit{3}
both show that the log-transformed observable is more linearly related to
$[\text{E}]$ across a wide dynamic range than $\Fk$ itself.

\subsection{Dynamic range}

The minimum detectable $\Fk$ is $\Fk^\text{min} = 3\sigma_{\Fk}(0)
\approx 20\%$ (Section~\ref{sec:Fk}). The maximum usable $\Fk$ is
approximately 0.95 (near-complete cleavage, beyond which the MB signal
is indistinguishable from noise). The corresponding apparent rate range is:
\begin{equation}
  \kk^\text{min} = \frac{-\ln(1-0.20)}{T} = \frac{0.223}{T},
  \qquad
  \kk^\text{max} = \frac{-\ln(1-0.95)}{T} = \frac{3.00}{T}
  \label{eq:dynrange}
\end{equation}
giving a dynamic range of $3.00/0.223 = 13.5\times$ in apparent rate, or
approximately 1.1 decades in $\kk$. In the linear (unsaturated) regime
where $\kk \approx \kcat[\text{E}]/(\Km + \Gtot(0))$, $\kk$ is proportional
to $[\text{E}]$, so 1.1 decades in $\kk$ corresponds to 1.1 decades in
$[\text{E}]$. The clinically relevant MMP-9 range is 6--30~nM
(chronic vs healing wound, Layer~0), spanning 0.7 decades. This fits
comfortably within the 1.1-decade window, confirming that the biosensor
can discriminate healing from chronic wounds across the full physiological
range with the stated noise floor.

\subsection{Inversion algorithm}

Given a measured $\Fk$ and its uncertainty $\sigma_{\Fk}$, the enzyme
concentration is extracted by solving Eq.~\eqref{eq:intMM} for $[\text{E}]$:
\begin{equation}
  [\text{E}] = \frac{1}{\kcat\,T}\left[\Km\ln\!\left(\frac{1}{1-\Fk}\right)
  + \Gtot(0)\,\Fk\right]
  \label{eq:inversion}
\end{equation}
This is a closed-form expression: given the kinetic parameters $\kcat$,
$\Km$, $\Gtot(0)$, and the measured $\Fk$, no numerical solver is needed.

The uncertainty on $[\text{E}]$ is:
\begin{equation}
  \sigma_{[\text{E}]} = \frac{1}{\kcat\,T}\left(
  \frac{\Km}{1-\Fk} + \Gtot(0)
  \right)\sigma_{\Fk}
  \label{eq:sigE}
\end{equation}
The prefactor $\Km/(1-\Fk) + \Gtot(0)$ diverges as $\Fk \to 1$, reflecting
the fact that near complete cleavage, a small absolute change in $\Fk$
corresponds to a large relative change in the number of remaining probes
and therefore to a large change in the inferred concentration. In practice,
readings with $\Fk > 0.90$ are flagged as "above range": $[\text{E}]$ is
reported as "$> [\text{E}]_\text{max}$" rather than a precise value with
a diverging uncertainty.

%=======================================================================
\section{Error Propagation Chain}
\label{sec:errprop}

The full error propagation from raw peak areas to enzyme concentration is:
\begin{equation}
  \underbrace{\sigma_{\AMB},\,\sigma_{\AFc}}_{\text{Layer~12}}
  \xrightarrow{\text{Eq.~\ref{eq:sigR}}}
  \sigma_R
  \xrightarrow{\text{Eq.~\ref{eq:sigrho}}}
  \sigma_\rho
  \xrightarrow{\text{Eq.~\ref{eq:sigFk}}}
  \sigma_{\Fk}
  \xrightarrow{\text{Eq.~\ref{eq:sigE}}}
  \sigma_{[\text{E}]}
  \label{eq:chain}
\end{equation}
Each arrow applies a specific formula derived above. The chain can be run
forward (from measured areas to final uncertainty) or in reverse (from
desired $\sigma_{[\text{E}]}$ to required per-scan precision $\sigma_{\AMB}$).
The reverse direction is used during instrument design to set the required
transimpedance amplifier noise floor.

Numerically, starting from $\sigma_{\AMB}/\AMB = 3\%$ per scan, the chain
gives (at $\Fk = 0.5$, $T = 30$~min, $N_\text{active} = 4$ electrodes averaged):
\begin{align*}
  \sigma_R/R &= 4.2\%\;\;(\text{single scan}) \\
  \sigma_\rho/\rho &= 4.8\%\;\;(\text{single electrode}) \\
  &= 2.4\%\;\;(\text{4 electrodes}) \\
  \sigma_{\Fk} &= (1-0.5) \times \sqrt{2} \times 2.4\% = 1.7\% \\
  \sigma_{[\text{E}]} &\approx \frac{\Km + \Gtot(0)}{\kcat T (1-0.5)}\times 1.7\%\;\;
  (\text{using Eq.~\ref{eq:sigE}})
\end{align*}
A 1.7\% absolute uncertainty on $\Fk = 0.5$ corresponds to approximately
a 3.4\% relative uncertainty on $\Fk$ at the midpoint, translating to
approximately 3--5\% relative uncertainty on $[\text{E}]$ (varying with the
exact slope of the calibration curve at that concentration).

%=======================================================================
\section{Electrode Layout and Averaging}
\label{sec:layout}

The sensor chip contains a fixed number of electrodes per patient cartridge.
The ratiometric pipeline aggregates readings from multiple electrodes of the
same type using a weighted mean:
\begin{equation}
  \bar{\kk} = \frac{\sum_i w_i\,\kk^{(i)}}{\sum_i w_i},
  \qquad
  w_i = \frac{1}{\sigma_{\kk^{(i)}}^2}
  \label{eq:wmean}
\end{equation}
where the weights are the inverse variances of the per-electrode apparent rates.
Electrodes flagged by Layer~12 quality diagnostics are assigned $w_i = 0$.
The uncertainty on $\bar{\kk}$ is:
\begin{equation}
  \sigma_{\bar{\kk}}^2 = \frac{1}{\sum_i w_i}
  \label{eq:sigkmean}
\end{equation}
The active-to-control electrode pairing is done before averaging: each
active electrode $i$ is paired with the spatially nearest control electrode
$j$ on the chip (to minimise differential fouling, which can have spatial
gradients across the chip). The double ratio $\Fk^{(i)}$ uses $\rho_\text{control}^{(j)}$
from its paired control, not the mean of all controls.

\subsection{Minimum electrode count and redundancy}

The minimum number of active electrodes needed to resolve $[\text{E}]$ at the
LOD ($\Fk = 0.20 \pm 3\sigma_{\Fk}$) at a 95\% confidence level is determined
by requiring:
\begin{equation}
  3\,\frac{\sigma_{\Fk}}{\sqrt{N_\text{active}}} < \Fk^\text{min}
  \quad \Rightarrow \quad
  N_\text{active} > \frac{9\,\sigma_{\Fk}^2}{(0.20)^2}
  = \frac{9 \times (0.068)^2}{0.04} \approx 1.0
  \label{eq:Nmin}
\end{equation}
A single active electrode is sufficient at the LOD (with $\sigma_{\Fk} = 6.8\%$
from Eq.~\ref{eq:sigFk}). In practice, $N_\text{active} \geq 3$ is required
for redundancy (in case one electrode fails the Layer~12 quality diagnostics)
and to achieve $\sigma_{[\text{E}]}/[\text{E}] < 5\%$ at mid-range.

%=======================================================================
\section{Limitations}
\label{sec:limits}

\begin{enumerate}[leftmargin=*]

\item \textbf{The kinetic parameters $\kcat$, $\Km$, and $\Gtot(0)$ are
  not measured per-cartridge.} Equation~\eqref{eq:inversion} requires
  these values as inputs. They are taken from Layer~9 (where $\kcat$ and
  $\Km$ come from bulk solution kinetics) and from Layer~11 (where
  $\Gtot(0)$ is estimated from the MB peak area before cleavage using the
  Layer~11 equation). Any batch-to-batch variation in aptamer synthesis
  yield, thiol coupling efficiency, or gold deposition that shifts
  $\Gtot(0)$ outside the range $[1 \times 10^{-12},\,5 \times 10^{-12}]$
  mol\,cm$^{-2}$ will bias the inversion. A per-cartridge measurement of
  $\Gtot(0)$ from $R(0)$ using the Layer~11 calibration is required for
  absolute accuracy.

\item \textbf{The double-ratio cancellation of fouling is exact only if
  fouling is uniform across the chip surface.} Baker~2006~\cit{3} and
  Lubin and Plaxco~\cit{2} note that plasma protein adsorption varies
  with SAM composition. If the Fc-thiol region and the aptamer-thiol
  region adsorb proteins at different rates, the Fc amplitude on the active
  and control electrodes will drift differentially, violating the
  approximation in Section~\ref{sec:rho}. The Fc amplitude monitoring
  diagnostic (Layer~12) detects but does not correct this failure mode.

\item \textbf{The inversion formula Eq.~\eqref{eq:inversion} assumes the
  integrated MM model, which uses surface $\Km$ estimated from bulk
  $K_m$ and an assumed effective diffusion layer depth.} Layer~9
  Section~IV establishes that $\Km$ spans two orders of magnitude depending
  on the assumed $\delta_\text{eff}$. If $\Km$ is uncertain by a factor of
  10, the inversion maps $\Fk$ to the wrong $[\text{E}]$ by approximately
  the same factor in the $\Km$-dominated regime ($\Fk \ll 1$). Calibration
  with known enzyme concentration samples eliminates this systematic error
  empirically; the model then correctly predicts the \emph{shape} of the
  calibration curve even if the absolute scale is set by the calibrant.

\item \textbf{The uncertainty formula Eq.~\eqref{eq:sigFk} assumes
  Gaussian-distributed fit residuals.} If the baseline is non-polynomial
  (blocked by the Layer~12 quality diagnostic) or if the peak is
  non-Gaussian (double-MB population), the residuals are non-Gaussian and
  the propagated $\sigma_{\Fk}$ underestimates the true uncertainty.

\item \textbf{The D-enantiomer control must be verified stereospecifically
  resistant to MMP-9.} The design assumes MMP-9 cannot cleave
  D-SGKIPRTLTA. This assumption holds for natural L-specific metalloproteinases
  but must be confirmed experimentally: add saturating MMP-9 to a chip
  bearing only D-enantiomer electrodes and confirm $\rho_\text{control}$
  does not drop by more than 5\%. Any residual activity on the D-enantiomer
  causes $\Fk$ from Eq.~\eqref{eq:Fk_simp} to underestimate $p_\text{E}$.

\end{enumerate}

%=======================================================================
\section*{References}

\noindent All papers in \texttt{electrochemical\_device/papers/layer\_13/}.

\begin{enumerate}[label={[\arabic*]},noitemsep,leftmargin=*]

\item White RJ, Plaxco KW (2010). Exploiting binding-induced changes in
  probe flexibility for the optimization of electrochemical biosensors.
  \textit{Anal Chem} 82:73--76. DOI:~10.1021/ac9028595.
  \textbf{white2010.pdf} -- ratiometric signal construction, cancellation of
  electrode area variation, and SWV frequency optimisation confirmed.

\item Lubin AA, Plaxco KW (2010). Folding-based electrochemical biosensors:
  the case for responsive nucleic acid architectures. \textit{Acc Chem Res}
  43:496--505. DOI:~10.1021/ar900165x.
  \textbf{lubin\_plaxco2010.pdf} -- structure-switching E-AB biosensor
  architecture, background cleavage correction by control electrode, fouling
  cancellation via double normalisation, and dynamic range characteristics
  confirmed.

\item Baker BR, Lai RY, Wood MS, Doctor EH, Heeger AJ, Plaxco KW (2006).
  An electronic, aptamer-based small-molecule sensor for the rapid,
  label-free detection of cocaine in adulterated samples and biological
  fluids. \textit{J Am Chem Soc} 128:3138--3139. DOI:~10.1021/ja056957p.
  \textbf{baker2006.pdf} -- on-chip reference electrode stability, SWV
  parameter validation, and dynamic range in biological fluid confirmed.

\item Lubin AA, Lai RY, Baker BR, Heeger AJ, Plaxco KW (2010).
  Electrochemical, structure-switching biosensors: quantitative analysis
  of the usable dynamic range. \textit{Anal Chem} 78:5671--5677.
  DOI:~10.1021/ac0601417.
  \textbf{lubin2010.pdf} -- per-electrode $R(0)$ baseline normalisation,
  electrode-to-electrode $R(0)$ variation, and Fc surface coverage stability
  during assay confirmed.

\end{enumerate}

\end{document}
